\chapter{The frozen batch: what is settled and what is not}
\label{ch:status}

This chapter is the experimental record. It replaces an earlier version that reported results
from runs at three different EM iteration budgets, three different replicate counts and two
values of $\corr$ --- numbers that could not go in one table, and several of which turned out to
be artefacts of those settings rather than of the model. Chapter~\ref{ch:corrections} records
what those were and how they were caught.

Everything below comes from a single configuration, defined once in
\codefile{experiments/frozen\_config.py} and imported by every experiment.

\section{The configuration}

\begin{center}\small
\begin{tabular}{@{}ll@{}}
\toprule
$\corr = 0.85$ & one value everywhere; the earlier $0.8$ runs are retired \\
$\nsites = 32$ & chain length, which is also the ambient dimension \\
$\ngrid = 401$, $\halfwidth = 8$ & read off the grid/domain sweep, \S\ref{sec:rec-rung0} \\
twelve noise levels & log-spaced on $[0.05, 3.0]$ \\
sixteen replicates & everywhere \\
$\nseq \in \{32,\dots,4096\}$ & seven doublings \\
$256$ held-out sequences & drawn from a seed disjoint from every training seed \\
EM stopping & relative $10^{-9}$ on the log-evidence, cap $400$; \textbf{never} on $\corr$ \\
ring & $T = 12$, so $\nsites = 24$ \\
\bottomrule
\end{tabular}
\end{center}

\begin{pitfall}
A frozen configuration is only as good as its reach, and this one has now failed that test twice.
Four experiments carried private replicate knobs that never read \texttt{n\_seeds}, and every one
of them reported \texttt{is\_frozen: true} while running at three, four or eight replicates
(\S\ref{sec:corr-replicates}). The same disease then turned up in the iteration budget:
\texttt{em\_max\_iters} and \texttt{em\_shape\_tol} are declared in the config and read by
\emph{nothing}, while \codefile{exp\_07} passed a literal $120$
(\S\ref{sec:corr-capacity}). Both were invisible to the provenance stamp, which compares the
config against itself. \codefile{tools/check\_replicates.py} now parses the settings dictionaries
for both classes of knob.

The general lesson is worth stating once: a frozen configuration is a claim about what the
experiments read, and only a reader-side check can test that claim. Declaring a value is not the
same as consuming it.
\end{pitfall}

\section{Settled: the numerical foundation}
\label{sec:rec-rung0}

\paragraph{The discretisation is controlled on two axes, separately.} Truncation --- mass outside
$[-\halfwidth,\halfwidth]$ discarded rather than transported --- responds to the domain and not
the resolution: the residual falls from $3\times10^{-4}$ at $\halfwidth=4$ to $1.4\times10^{-8}$
at $\halfwidth=8$, essentially independently of $\ngrid$. Quadrature error falls by a factor of
four per halving of $\gridh$, the predicted second order for the trapezoidal rule. Neither sweep
alone certifies the other, which is why both are reported.

\paragraph{The pipeline is verified where the answer is known.} At $\ngrid=401$, $\halfwidth=8$
the discretised recursion reproduces the Gaussian closed form to $9.2\times10^{-15}$ relative,
and agrees with brute-force enumeration of the discretised posterior --- summing over all
$\ngrid^{\nsites}$ configurations on a short chain, sharing no code with the message passing ---
to $1.6\times10^{-14}$ on posterior means and $1.8\times10^{-15}$ on the log-evidence. Two
independent routes, not one self-consistent one.

\paragraph{What moment closure costs.} By Proposition~\ref{prop:closure-lmmse} the gap between
moment-projected and full sum-product is exactly the price of replacing the prior by its
covariance-matched Gaussian. Measured on the Laplace chain, the median relative score error falls
from $0.241$ at $t=0.05$ to $7.3\times10^{-6}$ at $t=3$ --- a factor of $3.3\times10^{4}$ across
the schedule, as noising Gaussianises the posterior and the innovation's shape stops mattering.

\section{Settled: learning the kernel}

\paragraph{Both optimisation routes reach the same place.} On the Gaussian kernel, gradient
ascent via Fisher's identity reaches the exact M-step's optimum to $0.021$ nats at its best step
size, with slightly better parameter errors ($0.0040$ against $0.0053$ on $\corr$). The
difference is robustness, not accuracy: EM's monotonicity violation is $0.0$ exactly and it needs
no step size, while gradient ascent at $\eta=2.0$ diverges ($\corr$ error $0.41$, monotonicity
violation $1641$) and at $\eta=8.0$ is unstable. That is the whole content of the comparison, and
it is why every rung above uses EM.

\paragraph{The gradient is exact.} Fisher's identity is verified against finite differences of
the discretised evidence to ${\sim}10^{-9}$. This check matters more than it looks: a fixed-point
residual near zero is evidence of a stationary point only if the derivative defining it is right,
and a sign error in exactly that derivative once survived every internal diagnostic
(\S\ref{sec:corr-gradient}).

\section{Settled: what the structure buys}
\label{sec:rec-efficiency}

Relative denoising error against grid BP under the true kernel, averaged over the noise schedule,
with the network's parameterisation chosen on a disjoint validation split:

\begin{center}\small
\begin{tabular}{@{}rccc@{}}
\toprule
$\nseq$ & network & EM--BP ($12$ free) & ratio \\
\midrule
32   & $0.6613 \pm 0.0026$ & $\mathbf{0.0703 \pm 0.0051}$ & $9.4$ \\
64   & $0.5896 \pm 0.0029$ & $\mathbf{0.0541 \pm 0.0035}$ & $10.9$ \\
128  & $0.4942 \pm 0.0015$ & $\mathbf{0.0361 \pm 0.0024}$ & $13.7$ \\
256  & $0.3473 \pm 0.0012$ & $\mathbf{0.0259 \pm 0.0014}$ & $13.4$ \\
512  & $0.2359 \pm 0.0010$ & $\mathbf{0.0233 \pm 0.0010}$ & $10.1$ \\
1024 & $0.1685 \pm 0.0006$ & $\mathbf{0.0185 \pm 0.0011}$ & $9.1$ \\
2048 & $0.1364 \pm 0.0007$ & $\mathbf{0.0172 \pm 0.0008}$ & $7.9$ \\
\bottomrule
\end{tabular}
\end{center}

All sixteen seeds, all seven doublings, $\pm$ one standard error over seeds. The
validation-selected parameterisation agrees with the test-set optimum in $99.4\%$ of the $1{,}344$
cells, so the oracle it replaces is worth nothing measurable. The network at $\nseq=2048$
($0.136$) has not reached the error the estimator attains at $\nseq=32$ ($0.070$); no crossing is
observed anywhere in the range, so no data-equivalence factor is quoted.

The ratio is not monotone --- it rises to $13.7$ at $\nseq=128$ and falls to $7.9$ at $2048$ ---
and that shape should not be over-read. Both arms improve with $\nseq$; the estimator starts so
far ahead that the ratio is a quotient of two moving quantities, and nothing here identifies a
mechanism for the turn. What the table supports is the statement that the margin is an order of
magnitude across five doublings, not a claim about how it scales.

\begin{pitfall}
These sixteen seeds ran at $120$ EM iterations, because \codefile{exp\_07} hard-coded that budget
while \texttt{FROZEN.em\_max\_iters} said $400$ and nothing read it (\S\ref{sec:corr-capacity}).
The rerun at the configured budget is Bocconi array \texttt{628943}.

\textbf{The obvious prediction about that rerun was wrong, and the way it was wrong is the
interesting part.} The reasoning was: the shape settles at a median of $229$ updates, so a run
stopped at $120$ is short of convergence, so the estimator is handicapped and the ratios are a
lower bound. On the first six seeds the direction \emph{reverses with sample size}. Held-out score
error of the estimator, same six seeds, $120$ against $400$ iterations:

\begin{center}\small
\begin{tabular}{@{}rccc@{}}
\toprule
$\nseq$ & EM at $120$ & EM at $400$ & \\
\midrule
32   & $0.0585$ & $0.0673$ & worse by $15\%$ \\
64   & $0.0501$ & $0.0626$ & worse by $25\%$ \\
256  & $0.0283$ & $0.0295$ & worse by $4\%$ \\
512  & $0.0223$ & $0.0200$ & better by $10\%$ \\
2048 & $0.0158$ & $0.0125$ & better by $21\%$ \\
\bottomrule
\end{tabular}
\end{center}

\noindent
Running the optimiser \emph{longer} makes the estimator \emph{worse} at small $\nseq$ and better
at large $\nseq$, so no fixed budget is neutral between the two arms of the comparison.

\paragraph{The mechanism, measured rather than inferred.}
\label{sec:em-overfitting}
An instrumented run
(\codefile{experiments/exp\_29\_em\_overfitting.py}) records the training marginal log-evidence and
the held-out score error on the \emph{same} iterates, which is what separates overfitting from a
bad optimum: overfitting means the training objective keeps improving while generalisation turns
around, whereas a bad optimum would show both stalling.

\begin{center}\small
\begin{tabular}{@{}rccc@{}}
\toprule
$\nseq$ & training $\Lik$ monotone $\uparrow$ & held-out best at & cost of running to $400$ \\
\midrule
32   & yes & iteration $130$ & $+2.8\%$ \\
64   & yes & iteration $40$  & $+46.0\%$ \\
2048 & yes & iteration $250$ & $+8.6\%$ \\
\bottomrule
\end{tabular}
\end{center}

\noindent
The training log-evidence increases monotonically to the last iterate at every size --- as EM
guarantees it must --- while the held-out error has an \emph{interior minimum} at every size. That
is overfitting, and the diagnosis is now a measurement rather than a plausible story.

Note what the third row says: the interior optimum is at $250$ even at $\nseq=2048$. The effect is
\textbf{not} confined to small samples, and the earlier reading --- ``$120$ wins below the
crossover, $400$ wins above'' --- was an artefact of comparing two arbitrary points either side of
an optimum that moves. A cap of $400$ overshoots at every size tested.

\paragraph{Two consequences.} First, ``run it to convergence'' is not automatically the honest
choice. The iteration count is a hyperparameter with a bias--variance trade-off like any other, and
fixing it a priori quietly favours one arm; picking it on the test set would be exactly the sin
\S\ref{sec:corr-flat} records. It is now chosen on the validation bundle, per noise level,
symmetrically with the network's parameterisation, which is what \texttt{629091} runs.

Second --- and this is why it is in this chapter rather than a footnote --- this is the
memorisation--generalisation question of \S\ref{sec:memorisation} appearing in a model with twelve
free parameters and no network anywhere. The training objective and the generalisation objective
come apart, early stopping is doing implicit regularisation, and all of it is visible in a setting
where the target score is computable in closed form. That is precisely what this setting was built
for, and it is the most promising thing found in this batch.

\paragraph{Still not claimed.} That the optimal stopping point scales with $\nseq$ in any
particular way. The three sizes here give $130$, $40$ and $250$, which is not monotone, and three
points chosen for a diagnostic are not a curve. Reading a scaling law off them would repeat
\S\ref{sec:corr-flat} exactly.
\end{pitfall}

\section{Settled: structure the marginals cannot see}

Chapter~\ref{ch:ring} in full. The result in one line: across $1{,}344$ cells the per-frame arm's
estimate of $p(\psi)$ moves $0.0000$ in total variation from its uniform initialisation, and
across the $480$ cells of the single-rotation variant its profile likelihood in $\psi$ has range
$0.00\times10^{0}$ nats --- a machine-precision zero, which is what
Theorem~\ref{thm:ring-blind} requires. The joint arm recovers the rotation to $0.15^\circ$ at
$t=0.05$, degrading to $7.09^\circ$ at $t=1.43$ and returning to the null by $t=3$.

\section{Settled: where the Markov assumption fails}
\label{sec:rec-nonmarkov}

Two contamination mechanisms, both with an exact reference, from the complete ten-task run
\texttt{627165}. \texttt{ratio\_to\_em} is the baseline's relative score error over EM--BP's, so
$>1$ means the estimator wins.

\begin{center}\small
\begin{tabular}{@{}lcccccc@{}}
\toprule
& \multicolumn{5}{c}{global latent, rank-one, strength $\beta$} \\
\cmidrule(lr){2-6}
& $0.00$ & $0.10$ & $0.25$ & $0.50$ & $1.00$ \\
\midrule
Gaussian, vs CNN & $18.3$ & $17.2$ & $6.99$ & $3.36$ & $2.22$ \\
Laplace, vs CNN  & $9.90$ & $10.6$ & $6.50$ & $3.05$ & $2.12$ \\
\midrule
& \multicolumn{5}{c}{long-range precision coupling, strength $\gamma$} \\
\cmidrule(lr){2-6}
& $0.00$ & $0.05$ & $0.10$ & $0.20$ & $0.40$ \\
\midrule
Gaussian, vs CNN & $24.1$ & $1.19$ & $\mathbf{0.98}$ & $0.86$ & $0.77$ \\
Gaussian, vs MLP & $37.0$ & $1.79$ & $1.23$ & $0.96$ & $0.83$ \\
\bottomrule
\end{tabular}
\end{center}

The estimator survives rank-one contamination essentially intact --- still winning at $\beta=1$,
where half the marginal variance is a shared constant --- and \textbf{inverts at
$\gamma\approx0.10$} under long-range coupling. The mechanism is exactly what the structure
predicts: a global latent makes the score residual rank one, which a chain absorbs; long-range
coupling is not low-rank and cannot be represented. This is the honest scope statement, and it is
backed by a complete run.

\begin{pitfall}
These runs predate \codefile{frozen\_config.py} and are \emph{not} at the frozen settings. They
are reported here, in the development document, and are deliberately not quoted in the note or
the workshop paper. The qualitative conclusion --- rank-one survivable, long-range not --- is
robust to the settings; the specific ratios are not to be pooled with
\S\ref{sec:rec-efficiency}.
\end{pitfall}

\section{In progress}
\label{sec:pending}

Four of the five reruns listed in the previous version have landed and moved into the sections
above; what follows is what they said and what is still out.

\paragraph{Landed: the parametric rate} (\texttt{628576}). At four replicates the fitted
$\log$--$\log$ slopes were $-0.196$ for $\corr$ and $-0.700$ for $\ivar$ against a predicted
$-0.500$, with a non-monotone $\corr$ curve --- not a confirmation of $\nseq^{-1/2}$, and the
claim was pulled from the note rather than caveated. At sixteen the variance slope is $-0.519$
($r=-0.99$, monotone), which is the rate; $\corr$ is monotone but shallower at $-0.273$. Half the
claim is restored, half is not, and the note says so rather than averaging them into a verdict.

\paragraph{Landed: clean against channel} (\texttt{628600}). Two slownesses, and they are
independent. \emph{Optimisation}: coordinates of one kernel settle at very different rates ---
median $80$ updates for $\corr$, $68$ for $\ivar$, $229$ for the excess kurtosis
(\texttt{628601}, below). \emph{Statistics}: fitting through the channel costs rate, not merely a
constant.

\begin{center}\small
\begin{tabular}{@{}lcc@{}}
\toprule
$\log$--$\log$ slope, sixteen replicates & $\corr$ & $\ivar$ \\
\midrule
clean transition pairs & $-0.52$ & $-0.50$ \\
through the channel    & $-0.44$ & $-0.36$ \\
\bottomrule
\end{tabular}
\end{center}

\noindent
The clean arm sits on the parametric rate on both coordinates, which is the check that the
estimator is not itself the bottleneck. It is also flatly incompatible with the withdrawn
discretisation-floor reading of \S\ref{sec:corr-flat}: the clean arm falls by a factor of $3.4$ on
$\corr$ and $4.3$ on $\ivar$ across the range, so there is no floor to see.

\paragraph{Landed: shape convergence at sixteen seeds} (\texttt{628601}). The settling table above.
Zero monotonicity violations across all sixteen traces. The kurtosis figure has a long tail --- up
to $638$ updates in the slowest seed --- which is what makes a $400$ cap a decision rather than a
formality.

\paragraph{Landed: sample efficiency at sixteen seeds} (\texttt{628571}). \S\ref{sec:rec-efficiency},
with the iteration-budget caveat recorded there.

\paragraph{Landed, partly: the confining potential} (Rung 4a, \texttt{629193}).
\label{sec:rung4a-status}
Three attempts, and the third is still running, so this is worth setting out in order.

\texttt{628434\_5} was killed by the six-hour wall having written nothing --- $960$ gradient fits at
${\sim}0.2$\,s per step is roughly ten hours, and the experiment writes once at the end. Sharding by
seed fixed that. \texttt{628942} then completed and produced nothing usable, for a reason that had
nothing to do with the queue: \texttt{part\_potential} generated two species at $\pm\omega$ while
\texttt{fit\_potential} scores every trajectory at the single scalar $\psi$ it is handed, so half the
data was evaluated under the wrong rotation. The estimator's escape was to collapse the ring. At
$r_*\to0$ the ring is rotation-invariant and so pays nothing for a wrong $\psi$; $r_*$ pinned to the
lower clip in $51.2\%$ of cells, and the sweep reported the \emph{marginal} arm beating the joint
arm, which is the reverse of what this rung exists to show.

\texttt{629193} (one species, $600$ steps, sixteen seeds, $960$ cells) repairs that:

\begin{center}\small
\begin{tabular}{lccc}
\toprule
 & $\lambda$ rel.\ err.\ (med.) & $r_*$ abs.\ err.\ (med.) & pinned at clip \\
\midrule
joint, two species (\texttt{628942}) & $1.844$ & $0.950$ & $51.2\%$ \\
joint, one species (\texttt{629193}) & $0.229$ & $0.020$ & $4.8\%$ \\
marginal, one species                & $0.429$ & $0.039$ & $4.4\%$ \\
\bottomrule
\end{tabular}
\end{center}

\noindent The joint arm's error falls monotonically with sample size --- $0.618$, $0.336$, $0.222$,
$0.181$, $0.118$ at $n = 128$ to $4096$ --- and that trend is reported here because it survives the
caveat below.

\begin{pitfall}
The arm \emph{comparison} from \texttt{629193} is not reported, because the two arms were not held to
the same convergence standard. At $t\le1$ the marginal arm hit the $600$-step cap in $95\%$ of cells
against the joint arm's $57\%$, so ``joint beats marginal'' partly means ``joint converged and
marginal did not''. The obvious control does not rescue it: only $13$ of $400$ pairs have both arms
converged.

The cause was the stopping rule. \texttt{plateau\_tol} was an \emph{absolute} tolerance on an
objective that sums over sequences, hence eight times stricter at $n{=}4096$ than at $n{=}512$: a
convergence criterion varying along the very axis this rung reports its scaling on. It is now
per-sequence, pinned by
\texttt{tests/test\_ring.py::test\_stopping\_rule\_does\_not\_depend\_on\_sample\_size}, which
replicates one dataset fourfold and asserts the fit stops at the same step.
\end{pitfall}

\paragraph{But more optimisation is not more accuracy, and that is the deeper problem.}
\label{sec:rung4a-nonmonotone}
Twelve cells were refit to $2400$ steps with the plateau test disabled, to see what the budget was
costing. The first cell inspected showed exactly the expected picture --- at $n{=}512$,
$t{=}0.2216$, the joint arm is identical to eight significant figures from step $600$ onward while
the marginal arm's $\lambda$ error falls from $0.482$ to $0.372$ to $0.335$ --- and it would have
been easy to stop there and report that the joint arm was converged and the marginal arm cut short.
Across all twelve cells that is not what happens:

\begin{center}\small
\begin{tabular}{lccc c}
\toprule
cell & \multicolumn{3}{c}{$\lambda$ rel.\ err.\ at $600 / 1200 / 2400$} & log-lik.\ monotone \\
\midrule
$n{=}512$, $t{=}0.2216$, joint    & $0.078$ & $0.078$ & $0.078$ & yes \\
$n{=}512$, $t{=}0.4665$, joint    & $0.029$ & $0.057$ & $0.060$ & yes \\
$n{=}512$, $t{=}0.2216$, marginal & $0.177$ & $0.066$ & $0.207$ & yes \\
$n{=}2048$, $t{=}0.2216$, marginal& $0.265$ & $0.075$ & $0.149$ & yes \\
\bottomrule
\end{tabular}
\end{center}

\noindent The objective rises monotonically in all twelve cells --- backtracking guarantees it ---
while the parameter error is non-monotone in four, and in two joint-arm cells the estimate at $2400$
steps is \emph{worse} than at $600$. Ascending the likelihood harder does not monotonically improve
recovery, because the finite-sample maximiser is not at the truth. This is the same phenomenon as
the EM overfitting of \S\ref{sec:em-overfitting}, in a different estimator: the optimisation budget
is a regularisation knob, not a convergence detail.

So \texttt{629962} should be read for what it is. Holding both arms to one $n$-independent
criterion removes a defect that was certainly there; it does not deliver two converged estimators,
and the honest fix --- selecting each arm's stopping point on held-out data, as
\S\ref{sec:rec-efficiency} now does for EM-BP and the network --- is not implemented for this rung.
Any joint-versus-marginal margin from \texttt{629962} carries that caveat.

\paragraph{Still out.}
\begin{itemize}\itemsep3pt
\item \textbf{Rung 4a at a matched convergence budget} (\texttt{629962}), eight seeds split by
      sample size into sixteen shards, at a $3000$-step cap under the per-sequence tolerance. Until
      it lands, the joint-versus-marginal margin above stays out of every document.
\item \textbf{Sample efficiency with both budgets selected} (\texttt{629985}), which will replace
      the table in \S\ref{sec:rec-efficiency}. \texttt{629091} selected EM-BP's iteration count
      on validation while the network trained to a fixed $20{,}000$ steps --- the same fairness
      error that validation selection was introduced to remove, relocated rather than fixed.
      Training length is a bias/variance knob for both arms: EM's held-out error has an interior
      minimum at every sample size (\S\ref{sec:em-overfitting}), and SGD's does too. Selecting
      one budget while pinning the other favours the selected arm, and here it favoured us.
      The network's training length is now chosen on the same bundle, jointly with its
      parameterisation. Each arm has one selected quantity; the network selects over a
      two-dimensional grid and EM-BP over one.
\end{itemize}

\section{Withdrawn, and not replaced}

Recorded in full in Chapter~\ref{ch:corrections}; listed here so that nothing in this chapter is
read as still standing.

\begin{itemize}\itemsep3pt
\item \textbf{The capacity attribution.} That enlarging the innovation mixture buys generative
      fidelity. Roughly $96\%$ of the apparent effect was convergence rate at a fixed iteration
      budget. The rerun that would settle it lost all but one of its array tasks, so one cell
      survives and the claim stands withdrawn and unreplaced (\S\ref{sec:corr-capacity}).
\item \textbf{The discretisation floor.} That the clean-arm variance error was floored by the
      grid. It was under-replicated, not floored, and the grid contributes a constant
      $4.4\times10^{-4}$ --- an order of magnitude too small to produce the effect
      (\S\ref{sec:corr-flat}).
\item \textbf{Everything downstream of the superseded generation protocol}, including the
      four-family comparison at matched covariance (\S\ref{sec:corr-generation}).
\end{itemize}

\section{Not started}

The architectural question. Sum-product on a chain suggests a network that is narrow but as deep
as twice the chain, unlike the tree-structured case where the natural depth is the tree's
\citep{mei2024unet,garnierbrun2024transformers}. Which architectures exploit Markovianity, and at
what depth, is what this setting was built to measure, and is the obvious continuation. The
receptive-field sweep that would begin it exists (\codefile{exp\_12}) but selected its radius on
the evaluation set, so its numbers are exploratory and are not reported anywhere.
